\documentclass[11pt,a4paper]{article}

\usepackage[utf8]{inputenc}
\usepackage[T1]{fontenc}
\usepackage{amsmath,amssymb}
\usepackage{hyperref}
\usepackage[margin=2.5cm]{geometry}
\usepackage{graphicx}
\usepackage{natbib}

\title{Ertel Potential Vorticity Computation on Pressure Levels}
\author{PV\_Ertel\_Computation Project \\ \texttt{https://github.com/yanxingjianken/PV\_Ertel\_Computation}}
\date{v0.2.0 — June 2026}

\begin{document}
\maketitle

\section{Ertel PV in Height Coordinates}

Ertel's potential vorticity \citep{ertel1942} in height coordinates $(x,y,z)$ is:

\begin{equation}\label{eq:epv_z}
\mathrm{EPV} = \frac{1}{\rho}\,(\vec{\omega}_a \cdot \nabla\theta)
\end{equation}

where:
\begin{itemize}
  \item $\rho$ = density [kg\,m$^{-3}$]
  \item $\vec{\omega}_a = \nabla \times \vec{u} + 2\vec{\Omega}$ = absolute vorticity [s$^{-1}$]
  \item $\theta = T(p_0/p)^{R_d/c_p}$ = potential temperature [K]
\end{itemize}

In component form:
\begin{align}
\vec{\omega}_a &= \left(
  \frac{\partial w}{\partial y} - \frac{\partial v}{\partial z},\;
  \frac{\partial u}{\partial z} - \frac{\partial w}{\partial x},\;
  \frac{\partial v}{\partial x} - \frac{\partial u}{\partial y} + f
\right) \\
\mathrm{EPV} &= \frac{1}{\rho}\left[
  \left(\frac{\partial w}{\partial y} - \frac{\partial v}{\partial z}\right)\frac{\partial\theta}{\partial x}
+ \left(\frac{\partial u}{\partial z} - \frac{\partial w}{\partial x}\right)\frac{\partial\theta}{\partial y}
+ \left(\zeta + f\right)\frac{\partial\theta}{\partial z}
\right]
\end{align}

where $\zeta = \partial v/\partial x - \partial u/\partial y$ is the relative vorticity and $f = 2\Omega\sin\phi$ is the Coriolis parameter.

The MPAS-Atmosphere model \citep{skamarock2012} computes EPV on its native
height-based, terrain-following vertical coordinate using finite-volume flux
divergence on a Voronoi mesh (see \texttt{mpas\_pv\_diagnostics.F}).

\section{Isobaric Transformation}

Under the hydrostatic approximation $\partial p/\partial z = -\rho g$, the
vertical coordinate transforms from $z$ to $p$:

\begin{equation}
\frac{\partial}{\partial z} \rightarrow -\rho g\frac{\partial}{\partial p}
\end{equation}

Applying this to Eq.~\eqref{eq:epv_z}:

\begin{align}
\mathrm{EPV} &= \frac{1}{\rho}\bigg[
  \frac{\partial w}{\partial y}\frac{\partial\theta}{\partial x}
  - \frac{\partial w}{\partial x}\frac{\partial\theta}{\partial y}
  + \rho g\left(
    \frac{\partial v}{\partial p}\frac{\partial\theta}{\partial x}
    - \frac{\partial u}{\partial p}\frac{\partial\theta}{\partial y}
    - (\zeta + f)\frac{\partial\theta}{\partial p}
  \right)
\biggr]
\end{align}

Neglecting terms involving the horizontal gradient of vertical velocity $w$
(valid at synoptic and larger scales), we obtain the \textbf{isobaric Ertel PV}:

\begin{equation}\label{eq:epv_p}
\boxed{\mathrm{PV} = -g\left[
  \frac{\partial v}{\partial p}\frac{\partial\theta}{\partial x}
  - \frac{\partial u}{\partial p}\frac{\partial\theta}{\partial y}
  + (f + \zeta)\frac{\partial\theta}{\partial p}
\right]}
\end{equation}

In PVU ($1\;\mathrm{PVU} = 10^{-6}\,\mathrm{K\,m^2\,kg^{-1}\,s^{-1}}$):

\begin{equation}
\mathrm{PV}_{\mathrm{PVU}} = \mathrm{PV}_{\mathrm{SI}} \times 10^6
\end{equation}

\section{Full 3-Term vs Simplified Formula}

Eq.~\eqref{eq:epv_p} contains three terms:

\begin{enumerate}
  \item \textbf{Vortex stretching}: $(f + \zeta)\,\partial\theta/\partial p$
  \item \textbf{Horizontal shear (1)}: $(\partial v/\partial p)(\partial\theta/\partial x)$
  \item \textbf{Horizontal shear (2)}: $-(\partial u/\partial p)(\partial\theta/\partial y)$
\end{enumerate}

The simplified (ERA5-like) formula retains only term (1):

\begin{equation}\label{eq:simple}
\mathrm{PV}_{\mathrm{simple}} = -g\,(f + \zeta)\,\frac{\partial\theta}{\partial p}
\end{equation}

This approximation is valid when the horizontal shear of the vertical wind
shear is small compared to planetary vorticity stretching — true for most
synoptic-scale flows but not near strong fronts or jets.

In practice, Eq.~\eqref{eq:simple} is what ERA5 reports as potential vorticity
on pressure levels. Our validation against ERA5 native PV shows that the
simplified formula matches better (RMSE $\sim\!0.21$ PVU at 500\,hPa) than
the full 3-term formula (RMSE $\sim\!0.35$ PVU), because the additional
shear terms introduce numerical noise that outweighs their physical
contribution at the resolutions tested.

\section{Moisture Correction}

In the presence of water vapor, the virtual temperature $T_v$ replaces $T$:

\begin{equation}
T_v = T\,(1 + 0.61\,q)
\end{equation}

where $q$ is specific humidity [kg\,kg$^{-1}$]. The virtual potential
temperature is then:

\begin{equation}
\theta_v = T_v\left(\frac{p_0}{p}\right)^{R_d/c_p}
\end{equation}

Using $\theta_v$ instead of $\theta$ in Eq.~\eqref{eq:epv_p} provides a
moisture correction that is largest in the tropical boundary layer (where
$q \sim 0.015$--$0.020$\,kg\,kg$^{-1}$, giving $T_v - T \sim 3$--$4$\,K).
At mid-tropospheric levels and above, $q \ll 0.001$ and the correction is
negligible.

Our ERA5 validation shows that the moisture correction reduces RMSE by
${\sim}0.001$ PVU at 500\,hPa — a negligible improvement for synoptic-scale
studies, but potentially relevant for tropical or boundary-layer
applications.

\section{Numerical Discretization}

\subsection{Vertical Derivatives}\label{sec:vert}

Following MPAS \texttt{calc\_vertDeriv\_center} and
\texttt{calc\_vertDeriv\_one}, we use:

\begin{itemize}
  \item \textbf{Interior} ($k = 1,\ldots,N-2$): centred difference
    \begin{equation}
    \left.\frac{\partial f}{\partial p}\right|_k = \frac{1}{2}\left(
      \frac{f_{k+1} - f_k}{p_{k+1} - p_k}
    + \frac{f_k - f_{k-1}}{p_k - p_{k-1}}
    \right)
    \end{equation}
  \item \textbf{Top} ($k = N-1$, lowest pressure): one-sided backward
    \begin{equation}
    \left.\frac{\partial f}{\partial p}\right|_{N-1} =
      \frac{f_{N-1} - f_{N-2}}{p_{N-1} - p_{N-2}}
    \end{equation}
  \item \textbf{Surface} ($k = 0$, highest pressure): one-sided forward
    \begin{equation}
    \left.\frac{\partial f}{\partial p}\right|_0 =
      \frac{f_1 - f_0}{p_1 - p_0}
    \end{equation}
\end{itemize}

Pressure levels are ordered surface $\to$ top of atmosphere (highest to
lowest pressure) throughout.

\subsection{Horizontal Derivatives}

On a regular latitude--longitude grid:

\begin{align}
\frac{\partial f}{\partial y} &= \frac{1}{R_E}\frac{\partial f}{\partial\phi}
  \approx \frac{1}{R_E}\frac{\Delta f}{\Delta\phi} \\
\frac{\partial f}{\partial x} &= \frac{1}{R_E\cos\phi}\frac{\partial f}{\partial\lambda}
  \approx \frac{1}{R_E\cos\phi}\frac{\Delta f}{\Delta\lambda}
\end{align}

where $R_E = 6.371 \times 10^6$\,m and $\Delta\phi$, $\Delta\lambda$ are
computed via \texttt{numpy.gradient}.

\subsection{Pole Handling}

At the poles ($\phi = \pm 90^\circ$), $\cos\phi = 0$ and $dx \to 0$,
leading to a singularity in $\partial/\partial x$. We clip $\cos\phi$ to a
minimum value corresponding to half a grid spacing from the pole:

\begin{equation}
\cos\phi_{\min} = \cos(90^\circ - 0.5\,\Delta\phi)
\end{equation}

For a $1^\circ$ grid, $\cos\phi_{\min} \approx 0.0087$.

\section{Hybrid Sigma-Pressure Levels}

For models using a hybrid vertical coordinate (e.g., CESM2/CAM6), the
pressure at model level $k$ is:

\begin{equation}\label{eq:hybrid}
p_k = a_k\,p_0 + b_k\,p_s
\end{equation}

where $a_k$ (``hyam'') and $b_k$ (``hybm'') are the hybrid coefficients at
model midpoints, $p_0 = 100{,}000$\,Pa is the reference pressure, and $p_s$
is the surface pressure.

When $a_k$ and $b_k$ are available, the per-column actual pressure $p_k(i,j)$
replaces the 1-D approximate pressure level array in all vertical derivative
computations (Sec.~\ref{sec:vert}). This is critical for accurate PV near
the surface, where hybrid levels deviate most from pure pressure levels.

\paragraph{Note for CESM2-LENS2 users:} The static grid file on AWS S3
(\texttt{atm/static/grid.zarr}) contains NaN values for the hybrid
coefficients as of 2026-06-01. Users must obtain $a_k$ and $b_k$ from
alternative sources (e.g., CAM initial condition files, CESM source code)
to use this feature.

\section{Validation Methodology}

\subsection{ERA5 Comparison}

We validate against ERA5 native potential vorticity on pressure levels
(product: \texttt{reanalysis-era5-pressure-levels}, variable:
\texttt{potential\_vorticity}).

For a test case (2025-01-08 00Z, all 37 pressure levels, global $0.25^\circ$):

\begin{center}
\begin{tabular}{lccc}
\hline
Level & RMSE (simple, PVU) & Bias (PVU) & Correlation \\
\hline
850\,hPa & 3.53 & +0.10 & 0.28 \\
500\,hPa & 0.21 & +0.01 & 0.96 \\
250\,hPa & 1.09 & -0.03 & 0.97 \\
\hline
\end{tabular}
\end{center}

The poor 850\,hPa comparison is expected: pressure-level PV products are
least accurate near the surface where terrain-following coordinates differ
most from pressure coordinates. The excellent 500\,hPa and 250\,hPa
comparisons confirm the validity of the isobaric PV computation for the
free troposphere and lower stratosphere.

\subsection{CESM2-LENS2 Quick-Look}

Applied to member \texttt{r10i1181p1f1} (2010-01-01), with 32 hybrid levels
using approximate pressure. PV values are physically reasonable:
$-1.8$ to $+3.1$\,PVU at $\sim$525\,hPa, and
$-10.3$ to $+11.8$\,PVU at $\sim$233\,hPa.

\section*{References}

\bibliographystyle{plainnat}
\begin{thebibliography}{10}

\bibitem{ertel1942}
Ertel, H. (1942).
Ein neuer hydrodynamischer Wirbelsatz.
\textit{Meteorologische Zeitschrift}, 59, 277--281.

\bibitem{holton2013}
Holton, J.~R.\ \& Hakim, G.~J. (2013).
\textit{An Introduction to Dynamic Meteorology}, 5th ed.
Academic Press.

\bibitem{bluestein1992}
Bluestein, H.~B. (1992).
\textit{Synoptic-Dynamic Meteorology in Midlatitudes}, Vol.~I.
Oxford University Press.

\bibitem{skamarock2012}
Skamarock, W.~C.\ et al. (2012).
A description of the Advanced Research WRF version 3.
NCAR Technical Note NCAR/TN-475+STR.

\bibitem{mpas}
MPAS-Atmosphere source code.
\texttt{mpas\_toolchain/mpas/src/core\_atmosphere/diagnostics/mpas\_pv\_diagnostics.F}

\bibitem{era5}
Hersbach, H.\ et al. (2020).
The ERA5 global reanalysis.
\textit{Quarterly Journal of the Royal Meteorological Society}, 146, 1999--2049.

\bibitem{cesm2le}
Rodgers, K.~B.\ et al. (2021).
The CESM2 Large Ensemble.
\textit{Earth System Dynamics}, in preparation.
Data: \url{https://registry.opendata.aws/ncar-cesm2-lens/}

\end{thebibliography}

\end{document}
